Fix the convention $\mathbb N=\{1,2,3,\dots\}$ and $[1,N]=\{1,2,\dots,N\}$.  For $A\subseteq\mathbb N$ define its \emph{Schnirelmann density}
\[
 d_s(A):=\inf_{N\ge 1}\frac{|A\cap[1,N]|}{N}\in[0,1].
\]
For $b\in\mathbb Z$ write $A+b:=\{a+b:a\in A\}$.

Call $B\subseteq\mathbb N$ an \emph{additive basis} if there exists $k\in\mathbb N$ and $n_0$ such that every integer $n\ge n_0$ can be written as a sum of $k$ (not necessarily distinct) elements of $B$.

The question asks whether there exists a set $B\subseteq\mathbb N$ which is \emph{not} an additive basis, but for which there exists a function $f:(0,1)\to(0,\infty)$ such that for every $\alpha\in(0,1)$, every $A\subseteq\mathbb N$ with $d_s(A)=\alpha$, and every $N\ge 1$, there exists $b\in B$ (allowed to depend on $A$ and $N$) such that
\[
\bigl|(A\cup (A+b))\cap[1,N]\bigr|\ \ge\ (\alpha+f(\alpha))\,N.
\]
Edge cases: for $\alpha=0$ or $\alpha=1$ the requirement on $f(\alpha)$ is not imposed.
